% ---
% filename : ["electrotechnique_circuits_ac_20260428_251_circuit_rc_parallele.tex"]
% id : "electrotechnique_circuits_ac_20260428_251"
% title : "Circuit RC parallèle : grandeurs physiques"
% domain : "Électrotechnique"
% subdomain : "Circuits AC"
% tags : ["circuit rc", "parallèle", "impédance", "puissance active", "puissance réactive", "facteur de puissance", "ts"]
% difficulty : 2
% time_solve : 16  # Calcul : 2 (difficulté) * (4 questions * 2)
% has_octave : true
% octave_files : ["electrotechnique_circuits_ac_20260428_251_circuit_rc_parallele.m"]
% author : "om"
% ---
\documentclass[a4paper,11pt,answers,addpoints]{exam}
% Fichier : preambule_generer_exercice.sty
%
% ---------------------------------------------------------
% LANGUE ET ENCODAGE
% ---------------------------------------------------------
\usepackage[utf8]{inputenc}
\usepackage[T1]{fontenc}   % Ajout recommandé (manquait)
\usepackage[french]{babel} % Ajout pour la langue française

% ---------------------------------------------------------
% GÉOMÉTRIE ET MISE EN PAGE
% ---------------------------------------------------------
\usepackage[left=1.7cm,right=1.5cm,top=2cm,bottom=1cm]{geometry}
\usepackage{microtype}      % Meilleure répartition du texte
\usepackage{float}          % Positionnement [H]
\usepackage{needspace}      % Saut conditionnel intelligent

% Éviter les veuves/orphelines (optionnel, à décommenter si besoin)
% \widowpenalty=10000
% \clubpenalty=10000

% Espacement vertical flexible
\raggedbottom
\setlength{\parskip}{0.5em plus 0.2em minus 0.1em}
\setlength{\parindent}{0pt}

% ---------------------------------------------------------
% MATHÉMATIQUES
% ---------------------------------------------------------
\usepackage{amsmath, amssymb, bm, esvect}

% Vecteurs
\newcommand{\V}{\overrightarrow}
\def\vect#1{\overrightarrow{\strut#1}}
\providecommand{\Degrees}[0]{\ensuremath{^\circ}}

% ---------------------------------------------------------
% UNITÉS ET GRANDEURS
% ---------------------------------------------------------
\usepackage{siunitx}
\sisetup{output-decimal-marker={,}}
\DeclareSIUnit \var {var}
\DeclareSIUnit \VA {VA}
\DeclareSIUnit \kVA {kVA}
\DeclareSIUnit[quantity-product={}] \degree{\text{\textdegree}}

% ---------------------------------------------------------
% GRAPHIQUES (TikZ + CircuitikZ + PGFPlots)
% ---------------------------------------------------------
\usepackage{tikz}
\usetikzlibrary{
	arrows.meta,
	intersections,
	decorations.pathreplacing,
	calligraphy,
	positioning
}

\usepackage[european,straightvoltages,EFvoltages]{circuitikz}
\usepackage{tkz-fct}
\usepackage{pgfplots}
\pgfplotsset{compat=1.12}

% Symbole étoile-triangle personnalisé
\newcommand{\wye}{%
	\mathbin{\tikz[x=1ex,y=1ex]{%
			\draw[line width=.1ex] (0,0)--(45:1)--++(-45:1) (45:1)--++(0,1);
	}}%
}

% Sankey diagram
\usepackage{sankey}
\pgfdeclarelayer{background}
\pgfdeclarelayer{foreground}
\pgfdeclarelayer{sankeydebug}
\pgfsetlayers{background,main,foreground,sankeydebug}

% ---------------------------------------------------------
% TABLEAUX
% ---------------------------------------------------------
\usepackage{tabularx}
\usepackage{tabu}      % Alternative à tabularx (garde les deux, pas de redondance critique)
\usepackage{booktabs}  % Pour des tableaux propres

% ---------------------------------------------------------
% HYPERLIENS
% ---------------------------------------------------------
\usepackage[colorlinks=true,
menucolor=red,
breaklinks=true,
urlcolor=blue,
linkcolor=blue]{hyperref}

% ---------------------------------------------------------
% CODE (LISTINGS)
% ---------------------------------------------------------
\usepackage{xcolor}
\usepackage{listings}

\lstdefinestyle{octavestyle}{
	language=Matlab,
	basicstyle=\ttfamily\small,
	keywordstyle=\color{blue},
	commentstyle=\color{green!50!black},
	stringstyle=\color{red},
	stepnumber=1,
	frame=single,
	breaklines=true,
	breakatwhitespace=true,
	tabsize=4
}

% ---------------------------------------------------------
% MISE EN TEXTE DIVERSES
% ---------------------------------------------------------
\usepackage{enumitem}
\setlist[enumerate]{label=\Alph*)}
% ---------------------------------------------------------
% COMMANDES PERSONNELLES
% ---------------------------------------------------------
\newcommand\nomauteur{bdminasmorgul@protonmail.com}
\newcommand\dateTe{29.04.2026}
\newcommand\brancheTe{TS}

% ---------------------------------------------------------
% STYLE DE L'EXERCICE
% ---------------------------------------------------------
\pagestyle{headandfoot}
\firstpageheadrule
\runningheadrule

\firstpageheader{\brancheTe\ (\nomauteur)}{Élève : }{(\dateTe) \thepage/\numpages}
\runningheader{\brancheTe\ (\nomauteur)}{Élève : }{(\dateTe) \thepage/\numpages}

% Compteur de points en bas de page
\newcommand{\totalpointtts}{%
	\begin{tikzpicture}[remember picture, overlay]
		\node[anchor=south west, inner sep=0pt, outer sep=0pt]
		at ([xshift=0.5cm,yshift=0.5cm]current page.south west)
		{\rotatebox{90}{Total des points : \numpoints}};
	\end{tikzpicture}
}

\firstpagefooter{\totalpointtts}{}{}
\runningfooter{\totalpointtts}{}{}





% Options de l'exercice
\printanswers
\addpoints
\pointsinmargin
\bracketedpoints
\shadedsolutions
\unframedsolutions

% ---------------------------------------------------------
% DOCUMENT
% ---------------------------------------------------------
\renewcommand{\thepartno}{\arabic{partno}}
\renewcommand{\partlabel}{\thepartno)}
\begin{document}
\section*{Préparation TS PQ CFC ELMO, IELE, PELE}
	
\begin{questions}

\question Un circuit est composé d'un condensateur de capacité $C = \qty{6,37}{[\mu F]}$ monté en parallèle avec une résistance $R = \qty{500}{[\ohm]}$. L'ensemble est alimenté par une tension alternative de $U = \qty{500}{[V]}$ à une fréquence de $f = \qty{50}{[Hz]}$.

Déterminez l'ensemble des grandeurs physiques de ce circuit (courants, impédance, puissances et facteur de puissance).
\begin{parts}
	\part[3] Calcul de la réactance capacitive.
	\part[3] Calcul des intensités des courants dans les branches et du courant total.
	\part[3] Calcul de l'impédance totale.
	\part[3] Calcul du facteur de puissance de l'installation.
	\part[3] Calcul des puissances.
\end{parts}
\vspace{0.5em}
\needspace{50\baselineskip}
\begin{solution}
	\begin{parts}
		\part Calcul de la réactance capacitive $X_C$ :
		\begin{align*}
			\omega &= 2 \cdot \pi \cdot f = 2 \cdot \pi \cdot \qty{50}{[Hz]} = \qty{314,16}{[rad/s]} \\[6pt]
			X_C &= \dfrac{1}{C \cdot \omega} = \dfrac{1}{\qty{6,37e{-6}}{} \cdot \qty{314,16}{}} = \qty{500}{[\ohm]}
		\end{align*}
		
		\part Calcul des intensités des courants dans les branches et du courant total :
		\begin{align*}
			I_R &= \dfrac{U}{R} = \dfrac{\qty{500}{[V]}}{\qty{500}{[\ohm]}} = \qty{1}{[A]} \\[6pt]
			I_C &= \dfrac{U}{X_C} = \dfrac{\qty{500}{[V]}}{\qty{500}{[\ohm]}} = \qty{1}{[A]} \\[6pt]
			I_{tot} &= \sqrt{I_R^2 + I_C^2} = \sqrt{\qty{1}{[A]}^2 + \qty{1}{[A]}^2} = \qty{1,414}{[A]}
		\end{align*}
		
		\part Calcul de l'impédance totale $Z$ et du facteur de puissance $\cos \varphi$ :
		\begin{align*}
			Z &= \dfrac{U}{I_{tot}} = \dfrac{\qty{500}{[V]}}{\qty{1,414}{[A]}} = \qty{353,6}{[\ohm]} \\[6pt]
			\cos \varphi &= \dfrac{I_R}{I_{tot}} = \dfrac{\qty{1}{[A]}}{\qty{1,414}{[A]}} = \qty{0,707}{}
		\end{align*}
		
		\part Calcul des puissances active $P$, réactive $Q$ et apparente $S$ :
		\begin{align*}
			P &= U \cdot I_R = \qty{500}{[V]} \cdot \qty{1}{[A]} = \qty{500}{[W]} \\[6pt]
			Q &= U \cdot I_C = \qty{500}{[V]} \cdot \qty{1}{[A]} = \qty{500}{[var]} \\[6pt]
			S &= U \cdot I_{tot} = \qty{500}{[V]} \cdot \qty{1,414}{[A]} = \qty{707}{[VA]}
		\end{align*}
	\end{parts}
	\vspace{0.5em}
	\needspace{30\baselineskip}
	\begin{verbatim}
		% Télécharger OCTAVE depuis https://wiki.octave.org/Using_Octave et l'installer
		% Puis copier-coller le code dans OCTAVE
		
		% Données du problème
		U = 500;        % Tension [V]
		R = 500;        % Résistance [Ohm]
		C = 6.37e-6;    % Capacité [F]
		f = 50;         % Fréquence [Hz]
		
		% Calculs intermédiaires
		omega = 2 * pi * f;
		printf('Pulsation omega = %.2f rad/s\n', omega);
		
		Xc = 1 / (C * omega);
		printf('Réactance capacitive Xc = %.2f Ohm\n', Xc);
		
		% Intensités
		Ir = U / R;
		Ic = U / Xc;
		Itot = sqrt(Ir^2 + Ic^2);
		printf('Courant Ir = %.2f A, Ic = %.2f A, Itot = %.3f A\n', Ir, Ic, Itot);
		
		% Impédance et facteur de puissance
		Z = U / Itot;
		cos_phi = Ir / Itot;
		printf('Impédance Z = %.2f Ohm, cos(phi) = %.3f\n', Z, cos_phi);
		
		% Puissances
		P = U * Ir;
		Q = U * Ic;
		S = U * Itot;
		printf('Puissance P = %.2f W, Q = %.2f var, S = %.2f VA\n', P, Q, S);
	\end{verbatim}
\end{solution}

\end{questions}
\end{document}